\section{Eliminating quotient singularities}\label{quot:section}
\label{sec:quotient}

\begin{proposition}[Smoothness under a finite cover by an abelian variety]
\label{quot:smooth-base}
Let \(f:X\to B\) be the original projective Lagrangian fibration in
the hypotheses of Theorem~\ref{thm:main}.  Suppose that \(B\) has
quotient singularities and that there is a finite surjective morphism
\[
                         p:A\longrightarrow B
\]
from an abelian variety \(A\) of dimension \(n\).  Then \(B\) is smooth.
\end{proposition}

The two inputs have different roles.  The original fibration supplies
the direct-image identity and the Fano property of \(B\).  The abelian
variety supplies translation fields, and the Fano property will make
the ramification line bundle of the lifted cover ample.  No Galois
property of \(p\) is assumed.  The proof therefore retains all arrows
of its relation over the canonical stack, together with their labels.

\subsection{The canonical stack and its labelled relation}
\label{quot:stack-field}

Write \(\mathfrak B\) for the smooth canonical Deligne--Mumford stack
of \(B\), and \(c:\mathfrak B\to B\) for its coarse morphism.
Its local presentations are \([U/G]\), where \(U\) is smooth and the
finite group \(G\) acts faithfully and freely in codimension one.
These presentations can equivalently be obtained from the maximal
charts of Proposition~\ref{loc:maximal-charts}: their normalized overlaps
are now étale everywhere.  In particular, \(\mathfrak B\) is proper,
has trivial generic stabilizer, and agrees with \(B\) over \(B_{\rm reg}\).
All fiber products of stacks below are \(2\)-fiber products.  An arrow in
such a fiber product retains its isomorphism label, even when its two
coarse coordinates coincide.

\begin{lemma}[Representable lifts]\label{quot:lifts}
The maps \(f\) and \(p\) lift to representable morphisms
\[
          \widetilde f:X\longrightarrow\mathfrak B,
       \qquad \widetilde p:A\longrightarrow\mathfrak B,
       \qquad c\widetilde f=f,\quad c\widetilde p=p.
\]
The first is projective and equidimensional of relative dimension
\(n\), with connected fibers.  The second is finite, flat, and surjective.
\end{lemma}

\begin{proof}
Take a local Galois chart \(U\to U/G\) as above.  Its non-étale
locus on the coarse space has codimension at least two.
The inverse image of that locus has codimension at least two in \(X\),
by equidimensionality, and in \(A\), by finiteness.
Consequently, the full normalizations of the components of the
pullbacks of \(U\to U/G\) that dominate these smooth sources are
finite, by excellence, and étale over the sources, by purity of the branch locus
\cite[Tag 0BMB]{StacksProject}.  Their deck actions are torsor actions:
the torsor identities hold on the dense étale locus and extend uniquely
to the finite étale covers.  The torsors and their equivariant maps to
\(U\) are precisely the data defining the two lifts.
Uniqueness on normalized overlaps makes these data compatible and
gives the stated representable lifts.  The normalized-source package
of Proposition~\ref{loc:maximal-charts} gives the asserted properties
of \(\widetilde f\) on each chart: its source is smooth and symplectic,
its map to \(U\) is projective with Lagrangian equidimensional fibers,
and Stein factorization
gives connected fibers and direct image \(\cO_U\).

The pullback of \(\widetilde p\) to \(U\) is finite, since it is finite
over the ordinary finite base change \(A\times_B U\).  It is a morphism
between smooth \(n\)-dimensional varieties and is flat by miracle flatness
\cite[Tag 00R4]{StacksProject}.  Thus \(\widetilde p\) is finite flat, and it is
surjective because its coarse morphism is surjective.
\end{proof}

The lifted cover defines the finite relation
\begin{equation}\label{quot:relation}
 \mathcal R=A\times_{\mathfrak B}A
   \ \substack{\xrightarrow{\ t\ }\\[-.6ex]\xrightarrow[\ s\ ]{\ }}
 A,\qquad
 \delta:\mathcal R\longrightarrow A,\qquad
 \delta(\xi)=t(\xi)-s(\xi).
\end{equation}
An arrow with source \(x\) and target \(y\) is a triple
\((x,y,\eta)\), where
\(\eta:\widetilde p(x)\simeq\widetilde p(y)\) is its isomorphism label.
Both projections are finite flat.  The relation is generically
étale over \(A\), hence reduced: a nilpotent section would be zero
generically and would contradict torsion-freeness over the integral
base \(A\).  It is pure-dimensional of dimension \(n\), and every
irreducible component is finite and surjective over \(A\) under both
projections.
If \(\nu:Z\to\mathcal R\) is the normalization of a component followed
by inclusion, we also write \(s,t\) for the induced coordinate maps
on \(Z\), and put \(\delta_Z=\delta\circ\nu\).

The identity section \(i:A\to\mathcal R\) is a closed immersion.
The inertia stack \(I(\mathfrak B)\) parametrizes a point of
\(\mathfrak B\) together with one of its automorphisms.  Thus
\[
                   \delta^{-1}(0)
                       =A\times_{\mathfrak B}I(\mathfrak B).
\]
In this zero-difference fiber the image of \(i\) is both open and
closed: the identity section of the finite unramified inertia is
open, and separatedness makes it closed.  We call its complement
the nonidentity part.  In particular, a point in that complement is
an arrow \((x,x,h)\) with \(h\ne1\); its label is retained even though
its two coordinates agree.

Retain the ample line bundle \(H\) on \(B\) chosen in
Section~\ref{sec:setup}, and put \(P=p^*H\) on \(A\).  The two
coordinate pullbacks of \(P\) to \(\mathcal R\) are isomorphic.
For an abelian subvariety \(D\subset A\), define its polarized
complement by
\begin{equation}\label{quot:orthogonal}
 D^\perp
   =\ker\bigl(A\xrightarrow{\lambda_P}\widehat A
                      \longrightarrow\widehat D\bigr)^0,
 \qquad
 \lambda_P(a)=t_a^*P\otimes P^{-1}.
\end{equation}
The superscript \(0\) means identity component.  The nondegenerate
rational alternating form of the polarization gives
\(\dim D^\perp=n-\dim D\), \((D^\perp)^\perp=D\), and reversal of
inclusions under this operation.

Here is the maximality argument that will eliminate nontrivial stabilizers.
If a nontrivial stabilizer existed, surjectivity of \(\widetilde p\)
would give a nonidentity arrow of difference zero.  Choose, among all
relation components containing such an arrow, one with difference
image of largest dimension.  We first construct one translation field
\(v\) on \(A\).  For every nonidentity automorphism \(h\), the
\(h\)-fixed projection of \(d\widetilde p(v)\) will be nonzero at every
lift of the corresponding inertia point.  This will force each of
these difference images to be proper
in \(A\).  The equality of the two polarization pullbacks will then
make the chosen image a torsion coset \(\tau+D\).  It contains zero,
so \(\tau\in D\) and the image is the subgroup \(D\) itself.  Its
properness makes \(K=D^\perp\) positive-dimensional.

Finally, ramification will supply another relation component with
differences outside \(D\).  On the normalization of the chosen
component, a \(K\)-translation will move a point above nonidentity
inertia over a ramification divisor.  A whole divisor through that
translated point will still map onto \(D\).  We will
place that divisor and its marked point in one component of the space
of composable arrows.  Composition on this component preserves the
marked nonidentity arrow and produces a difference image strictly
larger than \(D\), contradicting the choice of dimension.  The next
calculation constructs the translation field needed for the first step.

\subsection{A translation field along nonidentity inertia}

\begin{lemma}[A \(K\)-theory identity]\label{quot:k-identity}
In \(K^0(\mathfrak B)\),
\begin{equation}\label{quot:direct-image}
       \widetilde f_![\cO_X]
       =\sum_{j=0}^{n}(-1)^j[\Omega^j_{\mathfrak B}]
       =\lambda_{-1}(\Omega^1_{\mathfrak B}).
\end{equation}
Moreover, for some integer \(m>0\),
\begin{equation}\label{quot:nilpotence}
  \bigl([T_{\mathfrak B}]-n[\cO_{\mathfrak B}]\bigr)^m
                \lambda_{-1}(\Omega^1_{\mathfrak B})=0 .
\end{equation}
\end{lemma}

\begin{proof}
Use the global isomorphisms
\(\phi_j:R^j f_*\cO_X\xrightarrow{\sim}\Omega_B^{[j]}\)
chosen once on the original base in Section~\ref{sec:setup}, with
\(\phi_0\) the unit
isomorphism.  We apply the chart comparison of
Proposition~\ref{str:direct-images}.  Choose a connected affine quotient
chart \(U\to U/G\) with \(U/G\to B\) an integral separated étale
neighborhood, and view \(U\to\mathfrak B\) as its atlas map.
Write \(f_U:X_U\to U\) for the normalized lift,
put \(U^0=U\times_B B_{\rm reg}\), and let \(\pi:U^0\to B_{\rm reg}\)
be the induced étale morphism.  The complement of \(U^0\) has
codimension at least two.  Over this open the lift is the ordinary
base change of the original fibration.  Ordinary flat base change for
\(\pi\), followed by the étale cotangent identification, gives
\[
 (R^j(f_U)_*\cO_{X_U})|_{U^0}
   \simeq\pi^*(R^j f_*\cO_X|_{B_{\rm reg}})
   \xrightarrow{\ \pi^*\phi_j\ }\Omega^j_{U^0}.
\]
This open includes the discriminant over the regular part of \(B\);
the flat map used here is \(\pi\), so no smoothness of the fibers is
required for the base-change step \cite[Tag 02KH]{StacksProject}.

For \(j>0\), Ou's local reflexivity theorem applies to \(f_U\):
Proposition~\ref{loc:maximal-charts} supplies its projective,
equidimensional map with connected fibers and smooth canonical-trivial
source, and \(U\) is smooth
\cite[Theorem~1.3]{Ou2019}.  Both \(R^j(f_U)_*\cO_{X_U}\) and
\(\Omega_U^j\) are therefore their reflexive extensions from \(U^0\).
The displayed isomorphism extends uniquely across the complement.
For \(j=0\) the same conclusion is the unit isomorphism from
\((f_U)_*\cO_{X_U}=\cO_U\).  On a chart overlap the extended maps
agree over the inverse image of \(B_{\rm reg}\), by transitivity of
flat base change from the same \(\phi_j\).  Every component of the
overlap meets this open, since it maps étale to a chart and cannot be
supported over its codimension-at-least-two complement.  Uniqueness
of reflexive extension makes the maps agree everywhere.  Thus they
glue to degreewise isomorphisms
\[
       R^j\widetilde f_*\cO_X\simeq\Omega^j_{\mathfrak B}
                          \qquad(0\le j\le n).
\]
For \(j>n\), the restriction to \(U^0\) is zero by
Proposition~\ref{prop:background}, and reflexivity gives zero on \(U\).
Thus no higher direct-image terms occur.

These isomorphisms first give \eqref{quot:direct-image} in coherent
\(K\)-theory.  In this setting coherent and vector-bundle \(K\)-theory
agree.  Indeed, \(T_{\mathfrak B}\) is faithful on every stabilizer by
linearization of finite group actions.  A finite direct sum
\(\mathcal G\) of bounded tensor powers of this bundle and its dual
contains every irreducible stabilizer representation occurring on a
finite system of quotient charts.  Exactness of finite-group
invariants makes the evaluation
\[
 c^*c_*(\mathcal G^\vee\otimes\mathcal F)\otimes\mathcal G
                          \longrightarrow\mathcal F
\]
surjective for every coherent \(\mathcal F\).  For sufficiently large
\(k\), the coherent sheaf
\(c_*(\mathcal G^\vee\otimes\mathcal F)\otimes H^k\) is globally
generated on \(B\).  Its generators yield a surjection from a finite
sum of \(\mathcal G\otimes c^*H^{-k}\) onto \(\mathcal F\).
Iterating and using regularity of the smooth stack gives a finite
locally free resolution.  This proves the claimed comparison and
allows the identity to be pulled back in vector-bundle \(K^0\), even
along an inertia morphism which is not flat.

We spell out the nilpotence needed for
\eqref{quot:nilpotence}.  If \(V\) is a vector bundle of rank \(r\) on
a smooth projective variety \(Y\), then \([V]-r\) is nilpotent in
\(K^0(Y)\).  Indeed, pull back to the full flag bundle of \(V\).
This pullback is injective on \(K^0\), since pushforward of its
structure sheaf is \(\cO_Y\), with no higher direct images.
The pulled-back class is a sum of classes \([L]-1\) for line bundles.
For a globally generated line bundle \(L\), a finite set of generating
sections and its exact Koszul complex show that
\((1-[L^{-1}])^N=0\) for some \(N\).
An arbitrary line bundle is a quotient of two sufficiently positive
globally generated line bundles.  Therefore its class minus one is
also nilpotent: it is an invertible multiple of a difference of two
commuting nilpotent elements.  A finite sum of commuting nilpotents is
nilpotent.  This proves the assertion.

Apply it on the original projective \(X\) to
\(\widetilde f^*T_{\mathfrak B}\).
The projection formula and \eqref{quot:direct-image} give
\eqref{quot:nilpotence}, with one exponent \(m\) before any inertia
component is chosen.
\end{proof}

\begin{lemma}[A translation field with no fixed-direction zero]
\label{quot:translation-field}
For each connected component \(S\) of the nonidentity inertia of
\(\mathfrak B\), let \(j:S\to\mathfrak B\) be its natural morphism,
let \(h\) denote the tautological finite-order automorphism, and set
\[
 d=\dim S,\qquad
 W=A\times_{\mathfrak B}S,\qquad
 a:W\to A,\qquad q:W\to S,\qquad E=q^*T_S .
\]
There is a translation-invariant vector field
\(v\in\operatorname{Lie}(A)\) such that, simultaneously for every such
\(S\), the section
\begin{equation}\label{quot:projected-section}
 s_{v,S}\in H^0(W,E),\qquad
 s_{v,S}=\operatorname{pr}_{h=1}\bigl(d\widetilde p(v)|_W\bigr),
\end{equation}
has no zero.
\end{lemma}

\begin{proof}
An inertia component is a smooth proper stack: on a quotient chart its
presentation is a fixed locus \(U^h\) modulo the appropriate centralizer.
The fixed eigensubbundle of \(j^*T_{\mathfrak B}\) is \(T_S\).
Also, \(W\) is a scheme, since it is finite over \(A\).
The morphism \(W\to S\) is finite flat, being a base change of
\(\widetilde p\).
Thus \(W\) is projective, pure-dimensional of dimension \(d\), and
Cohen--Macaulay.  These assertions hold even if \(W\) is nonreduced.

Restrict \eqref{quot:nilpotence} to \(S\) and take the character of the
tautological automorphism. This is the inertia eigenbundle character
used in To{\"e}n's Riemann--Roch framework
\cite[Section~3.3 and Lemma~4.6]{Toen1999}; only its explicitly described
ring homomorphism and moving-factor cancellation are needed here.  More explicitly, for a bundle \(V\) on
\(\mathfrak B\), decompose
\(j^*V=\bigoplus_{\zeta} V_\zeta\) into the \(h\)-eigenbundles and use
the ring homomorphism
\[
             [V]\longmapsto\sum_{\zeta}\zeta[V_\zeta]
                \quad\hbox{in }K^0(S)\otimes_{\Z}\C .
\]
Pull this identity to \(W\) and take its Chern character in
\(H^{\rm even}(W,\C)\).
The degree-zero term coming from
\([T_{\mathfrak B}]-n\) is
\(\operatorname{tr}(h|T_{\mathfrak B})-n\), which is nonzero.
Indeed, the tangent action is faithful and all eigenvalues have
absolute value one; their sum can be \(n\) only if every eigenvalue is
one.  This would make \(h\) the identity by finite-order linearization.
The character of each moving cotangent factor in the alternating
exterior sum has nonzero degree-zero term
\((1-\zeta)^{\operatorname{rk}V_\zeta}\), for \(\zeta\ne1\).
The eigenbundle ranks are constant on connected \(S\), so each of
these degree-zero terms is the same nonzero scalar on every connected
component of \(W\).  Each full class is therefore a unit: every
positive-degree element of the cohomology of the finite-dimensional
projective scheme \(W\) is nilpotent.  Cancelling them gives
\[
                 \operatorname{ch}
                   \bigl(\lambda_{-1}(q^*\Omega^1_S)\bigr)=0.
\]
The term of degree \(2d\) on the left is \(c_d^{\rm top}(E)\), since for
Chern roots \(x_i\) of \(E\) the expression is
\(\prod_i(1-e^{-x_i})\).  Hence
\begin{equation}\label{quot:fixed-chern-zero}
                              c_d^{\rm top}(E)=0
                              \quad\hbox{in }H^{2d}(W,\C).
\end{equation}
In particular, a nonempty nonidentity inertia component of dimension
zero is impossible, as in that case this says \(1=0\).

Consider the bundle map
\[
      \alpha:\operatorname{Lie}(A)\otimes\cO_W
           =a^*T_A\longrightarrow E
\]
used in \eqref{quot:projected-section}.
For every reduced irreducible subvariety \(Y\subset W\) of dimension
\(s\), this map has rank at least \(s\) at its generic point.
In fact, \(Y\to S\) is finite onto its image and is generically
separable.  On a smooth dense open of \(Y\), and after an étale chart
on \(S\), its differential consequently has rank \(s\).
The differential identity for the fiber square says
\[
                     \alpha\circ d(a|_Y)=d(q|_Y),
\]
because the latter differential has values in the fixed tangent
subbundle.  This proves the rank assertion.

It follows that the locus on \(W_{\rm red}\) where
\(\operatorname{rk}\alpha\le r\) has dimension at most \(r\):
apply the preceding assertion to each of its irreducible components.
Over the stratum of rank \(r\), the incidence variety
\[
 \{(w,u)\in W_{\rm red}\times\operatorname{Lie}(A):
                                             \alpha_w(u)=0\}
\]
has dimension at most \(r+(n-r)=n\).
There are finitely many rank strata.  Thus a general \(u\) gives a
section whose zero scheme has at most zero-dimensional support.
It is a finite scheme because \(W\) is proper.

At every such zero, the \(d\) local entries of the section are a
system of parameters in a \(d\)-dimensional Cohen--Macaulay local ring,
hence a regular sequence \cite[Tag 02JN]{StacksProject}.
Let \([W]_d=\sum_i\operatorname{length}(\cO_{W,\eta_i})[W_i]\)
be the scheme-theoretic fundamental cycle, where the \(W_i\) are the
irreducible components with generic points \(\eta_i\).  The algebraic
zero-section formula gives
\[
       c_d^{\rm CH}(E)\cap[W]_d=[Z(s_{u,S})]_0
                              \quad\hbox{in }\mathrm{CH}_0(W)
\]
\cite[Tag 0FA9]{StacksProject}.  The complex Borel--Moore cycle class commutes with
vector-bundle Chern classes and proper pushforward
\cite[Chapter~19]{Fulton1998}.  With the same generic multiplicities
in \(\operatorname{cl}_{\rm BM}([W]_d)\), this gives
\[
 \operatorname{length} Z(s_{u,S})
   =\deg\bigl(c_d^{\rm CH}(E)\cap[W]_d\bigr)
   =\bigl\langle c_d^{\rm top}(E),
                  \operatorname{cl}_{\rm BM}([W]_d)\bigr\rangle=0,
\]
where the last equality is \eqref{quot:fixed-chern-zero}.  Each
nonempty isolated zero has strictly positive length.
Thus there are no zeros.
The inertia has only finitely many connected components, so the
intersection of the corresponding nonempty open sets of choices of
\(u\) supplies the single field \(v\) in the statement.
\end{proof}

\subsection{Difference images and translations}\label{quot:differences}

We now use the field to control every relation component meeting the
nonidentity part of \(\delta^{-1}(0)\).  The first step rules out a
dominant difference image; the polarization will then describe each
remaining image and its translation symmetry.

\begin{lemma}[Dominant differences avoid nonidentity inertia]
\label{quot:nondominance}
If an irreducible component of \(\mathcal R\) contains a nonidentity
arrow with difference zero, its difference image is a proper
subvariety of \(A\).
\end{lemma}

\begin{proof}
Let \(Z\) be the normalization of such a component, let \(z_0\in Z\)
map to the given arrow, and suppose that \(\delta_Z:Z\to A\) is
dominant.  Take its Stein factorization
\[
                       Z\xrightarrow{\beta}T\xrightarrow{\gamma}A.
\]
Here \(T\) is normal, \(\gamma\) is finite and dominant, and \(\beta\) has connected
fibers.  The connected fiber through \(z_0\) lies entirely over
nonidentity arrows: its image lies in \(\delta^{-1}(0)\), where the
identity and nonidentity parts are open and closed.

Put \(t_0=\beta(z_0)\), and take the field \(v\) of
Lemma~\ref{quot:translation-field}.
A finite dominant map to the smooth \(n\)-fold \(A\) is locally open.
For completeness, a sufficiently small local branch of \(T\) at
\(t_0\) is finite over a neighborhood of \(0\); its analytic image is
closed and has dimension \(n\), so contains a neighborhood of \(0\).
We can therefore lift a sequence
\(\exp(\varepsilon_\nu v)\), with
\(0\ne\varepsilon_\nu\to0\), to points of \(T\) tending to \(t_0\).
Lift these points farther to \(Z\).
Properness gives a subsequence tending to a point
\(z_\infty\in\beta^{-1}(t_0)\).
Its arrow is nonidentity, and the arrows in the sequence have
coordinates
\[
                  (x_\nu,y_\nu),\qquad
                  y_\nu=x_\nu+\exp(\varepsilon_\nu v),
             \qquad x_\nu,y_\nu\longrightarrow x_\infty .
\]

Choose a local quotient chart at \(\widetilde p(x_\infty)\), a lift
\(\psi\) of \(\widetilde p\) on a small analytic neighborhood of
\(x_\infty\), and finite-group linearizing coordinates at
\(\psi(x_\infty)\).
After passing to a subsequence the labels are one fixed group element
\(h\ne1\), fixing \(\psi(x_\infty)\), and
\[
                               \psi(y_\nu)=h\psi(x_\nu).
\]
Projecting to the \(h\)-fixed coordinate space yields
\[
        \operatorname{pr}_{h=1}
                \bigl(\psi(x_\nu+\exp(\varepsilon_\nu v))-\psi(x_\nu)\bigr)
                                                        =0.
\]
Division by \(\varepsilon_\nu\) and passage to the limit give
\(\operatorname{pr}_{h=1}d\psi_{x_\infty}(v)=0\).
This is a zero of the section \eqref{quot:projected-section} on the
inertia component containing \(h\), contrary to
Lemma~\ref{quot:translation-field}.
\end{proof}

\begin{lemma}[Difference images and complementary translations]
\label{quot:difference-image}
Let \(Z\) be the normalization of an irreducible component of
\(\mathcal R\), and suppose that its difference image \(D_Z\) has
dimension \(r<n\).
Then
\[
                              D_Z=\tau+D
\]
for an abelian subvariety \(D\subset A\) of dimension \(r\) and a
torsion point \(\tau\).
If \(K=D^\perp\), the image of \(Z\) in \(A\times A\) is stable under
diagonal translation by \(K\).  This action lifts uniquely to \(Z\);
both coordinate morphisms are equivariant, and \(\delta_Z\) is
invariant.
\end{lemma}

\begin{proof}
Let \(U_Z\subset D_Z\) be a nonempty open over which every
irreducible fiber component of \(\delta_Z\) has dimension \(n-r\),
let \(a\in U_Z\), and let \(F\) be any one of those components,
with its reduced structure.  On \(F\) the second coordinate map is \(t_a\) composed
with the first.  The equality of the two pulled-back polarizations
therefore says
\begin{equation}\label{quot:character-constraint}
                     (s|_F)^*\lambda_P(a)\simeq\cO_F .
\end{equation}
Let \(H_F\subset A\) be the connected abelian subvariety generated
by differences of points of \(s(F)\).  Since \(s\) is finite,
\[
                             \dim H_F\ge n-r.
\]
Resolve \(F\) and choose a base point on its resolution.
Its map to \(A\) factors through a surjective homomorphism from its
Albanese variety onto \(H_F\), followed by a translation.
The dual homomorphism
\(\widehat H_F\to\operatorname{Pic}^0(\widetilde F)\)
has finite kernel.  Equation \eqref{quot:character-constraint}
thus implies that the restriction of \(\lambda_P(a)\) to \(H_F\)
is torsion.

For a fixed abelian subvariety \(H\subset A\), the set of points with
this property is a countable union of torsion translates of
\(H^\perp\).  To check the assertion including disconnected kernels,
the homomorphism
\[
                      A\longrightarrow\widehat H,\qquad
                      a\longmapsto\lambda_P(a)|_H
\]
is surjective and has \(H^\perp\) as the identity component of its kernel;
the induced map \(A/H^\perp\to\widehat H\) is an isogeny.
The inverse image of any torsion point is consequently a finite
union of cosets of \(H^\perp\) represented by torsion points of \(A\).
Here torsion points lift through a quotient of abelian varieties;
this follows, for example, by choosing an isogeny complement to
the kernel.

There are only countably many abelian subvarieties of \(A\), since
their rational first-homology subspaces determine them.
There are also only countably many torsion points.
The open \(U_Z\) is therefore covered by countably many
closed torsion cosets, each of dimension at most \(r\).
A complex irreducible variety cannot be covered by countably many
proper closed subvarieties: apply the Baire theorem in a small
analytic ball in its smooth locus, and then use irreducibility.
Hence one such coset contains \(D_Z\).
The dimension bound forces equality, proving
\(D_Z=\tau+D\) with \(\dim D=r\).

This argument also determines the subgroup generated by every
component over a dense set of very general points.  Let \(a\) range
over \(U_Z\) minus every coset in the preceding countable collection
which does not contain \(D_Z\).  This complement is dense by the same
Baire argument.  For \emph{each} fiber component \(F\) over \emph{each}
such \(a\), its character constraint supplies a coset of
\(H_F^\perp\) of dimension at most \(r\) containing \(a\).
That coset must contain \(D_Z\) and hence equal it.
Thus
\[
                         H_F^\perp=D,\qquad H_F=D^\perp=K .
\]
The image \(s(F)\) lies in one \(K\)-coset and has dimension
\(n-r=\dim K\), so it is the whole coset.
Consequently the relation over every such very general \(a\) is
stable under diagonal \(K\)-translation.  These fibers are dense in
\(Z\): every nonempty Zariski open in \(Z\) has image containing a
nonempty Zariski open in \(D_Z\), which meets this set.  Taking closures
therefore proves this invariance on the full image
\(\Gamma\subset A\times A\).

The finite map from the original relation component to \(\Gamma\)
is generically one-to-one: over \(B_{\rm reg}\) there is a unique
isomorphism label for a pair of points.  Each component dominates
\(A\), so this open is dense on it.  Thus \(Z\) is also the
normalization of \(\Gamma\).
The morphism \(K\times Z\to\Gamma\) induced by simultaneous translation
is dominant and has normal source.  It therefore lifts uniquely to
the normalization \(Z\).
The action law, equivariance of the coordinate maps, and invariance
of the difference follow by uniqueness of this lift, or by checking
on the dense normal isomorphism locus and extending.
\end{proof}

\subsection{Ramification and maximality}\label{quot:maximality}

We now begin the contradiction proving
Proposition~\ref{quot:smooth-base}.  Suppose that \(\mathfrak B\) has
a nontrivial stabilizer.  Surjectivity of \(\widetilde p\) gives a
nonidentity point of \(\delta^{-1}(0)\).  Among the irreducible
components of \(\mathcal R\) containing such a point, choose one,
denoted \(\mathcal R_0\), whose difference image has largest
dimension \(r\).  Write
\[
       \nu:Z\longrightarrow\mathcal R_0\hookrightarrow\mathcal R
\]
for its normalization followed by inclusion.  By
Lemma~\ref{quot:nondominance}, \(r<n\).  Lemma~\ref{quot:difference-image}
makes \(\delta_Z(Z)\) a torsion coset; since it contains zero, it is
an abelian subvariety \(D\subsetneq A\).  Put \(K=D^\perp\), and fix
one point \(z_*\in Z\) above a nonidentity zero-difference arrow.

The next two lemmas provide exactly the data needed for one
composition.  The first gives a ramification divisor \(L\) and an
actual residual component containing all of \(i(L)\), with a
difference outside \(D\).  The second moves \(z_*\) over \(L\) and
places the translated point on a whole divisor whose difference
image is \(D\).

\begin{lemma}[Ramification and a residual relation]
\label{quot:ramification}
Let \(D\subsetneq A\) be an abelian subvariety and \(K=D^\perp\).
There is a prime component \(L\) of the ramification divisor of
\(\widetilde p:A\to\mathfrak B\) which maps surjectively to \(A/K\).
For such an \(L\), there is one irreducible component
\(\mathcal R'\ne i(A)\) of the actual relation \(\mathcal R\) such that
\[
                       i(L)\subset\mathcal R',
              \qquad   \delta(\mathcal R')\not\subset D .
\]
\end{lemma}

\begin{proof}
Since \(\widetilde p\) is generically étale, the determinant of
\(d\widetilde p\) defines an effective Cartier
ramification divisor \(\operatorname{Ram}(\widetilde p)\) on \(A\), and
\[
             \cO_A(\operatorname{Ram}(\widetilde p))
                        \simeq \widetilde p^*K_{\mathfrak B}^{-1}.
\]
There is no divisorial stabilizer on \(\mathfrak B\).
A sufficiently divisible power of its canonical bundle is therefore
the pullback of \(\cO_B(mK_B)\).  Thus
\[
 (\widetilde p^*K_{\mathfrak B}^{-1})^{\otimes m}
                           \simeq p^*\cO_B(-mK_B).
\]
The right side is ample because \(p\) is finite and \(-K_B\) is ample.
Hence the line bundle of the total ramification divisor is ample.

Since \(D\ne A\), the group \(K\) is positive-dimensional.
If every prime component of \(\operatorname{Ram}(\widetilde p)\) had proper
image in \(A/K\), a general \(K\)-coset would miss the entire divisor.
Its defining section would trivialize the restricted ample line
bundle on that positive-dimensional projective coset, a contradiction.
Thus one component \(L\) dominates \(A/K\); properness makes this
map surjective.  This uses ampleness of the \emph{total} ramification
divisor and does not assert that each of its components is ample.

The generic point of \(L\) lies over \(B_{\rm reg}\), because
finiteness preserves the codimension-at-least-two bound for the
inverse image of \(B_{\rm sing}\).
Near a general point of \(L\), the map between smooth spaces has
the transverse form
\[
               (u_1,u_2,\ldots,u_n)
                  \longmapsto(u_1^e,u_2,\ldots,u_n),
                         \qquad e>1.
\]
Its self fiber product has, besides the identity branch, branches
\(v_1=\zeta u_1,\ v_j=u_j\) for \(\zeta^e=1,\ \zeta\ne1\).
They meet the identity along \(L\).
There are only finitely many global irreducible components of
\(\mathcal R\).  Their closed intersections with \(i(L)\), excluding
the identity component, contain a dense open of \(i(L)\).
Irreducibility therefore gives a single residual component
\(\mathcal R'\ne i(A)\) containing all of \(i(L)\), including its
special points.

Suppose, for a contradiction, that \(\delta(\mathcal R')\subset D\).
Lemma~\ref{quot:difference-image} says that this difference image is a
torsion coset.  It contains zero, since \(i(L)\subset\mathcal R'\),
so it is an abelian subvariety \(D'\subset D\).
Its polarized complement \(K'=(D')^\perp\) contains \(K\).
The same lemma gives a \(K\)-action on the normalization \(Z'\) of
\(\mathcal R'\).

Consider the closed set
\[
                      F=\{x\in A:i(x)\in\mathcal R'\}.
\]
It is proper: otherwise \(\mathcal R'\) would contain \(i(A)\),
and the two irreducible \(n\)-dimensional components would be equal.
For any \(x\in F\), choose an individual preimage of \(i(x)\) in \(Z'\).
Its connected \(K\)-orbit remains over difference zero and remains
in the inverse image of the identity part of \(\delta^{-1}(0)\),
since that part is open and closed.  Its image at \(k\) is therefore
precisely \(i(x+k)\).  Hence \(F\) is \(K\)-invariant.
This reasoning does not require a global lift of \(i(L)\) to \(Z'\).

The divisor \(L\subset F\) is an irreducible component of \(F\),
because \(F\) is proper in the \(n\)-fold \(A\).
A connected group preserves each irreducible component of a closed
invariant set with finitely many components.  Thus \(K+L=L\).
But \(L\to A/K\) is surjective, which gives \(K+L=A\).
This contradiction proves the required assertion about
\(\delta(\mathcal R')\).
\end{proof}

\begin{lemma}[A divisor through the marked arrow]
\label{quot:marked-divisor}
Let \(\nu:Z\to\mathcal R\) be the normalization of an irreducible
component followed by inclusion, with difference image an abelian
subvariety \(D\subsetneq A\), and put
\(K=D^\perp\).
Suppose \(z_0\in Z\) lies over a nonidentity arrow of difference zero.
For \(L\) as in Lemma~\ref{quot:ramification}, there are \(k\in K\),
\(e_0=k\cdot z_0\), and a reduced prime divisor \(E\subset Z\)
through \(e_0\), such that \(x_0:=s(e_0)=t(e_0)\in L\),
\(\nu(e_0)\) is nonidentity, and
\[
                         s(E)=L,\qquad \delta_Z(E)=D .
\]
\end{lemma}

\begin{proof}
The action of \(K\) supplied by
Lemma~\ref{quot:difference-image} preserves difference zero.
Connectedness and the open-and-closed identity partition preserve
the nonidentity condition.  Since \(L\to A/K\) is surjective, a
translate of the first coordinate of \(z_0\) belongs to \(L\).
Choose such a \(k\), and put \(e_0=k\cdot z_0\) and
\(x_0=s(e_0)=t(e_0)\).  This is the translated normalization point
retained below.

The divisor \(L\) is Cartier because \(A\) is smooth.
Its local equation pulls back to a nonzero nonunit in the integral
local ring \(\cO_{Z,e_0}\).
A minimal prime over that equation has height one, by the principal
ideal theorem.  Taking its closure with reduced structure gives a
prime divisor \(E\) through \(e_0\) and contained in \(s^{-1}(L)\).
The finite map \(s:Z\to A\) preserves dimension, so its closed image
\(s(E)\) has dimension \(n-1\) and equals the irreducible divisor \(L\).

The image of the proper action morphism \(K\times E\to Z\) is closed
and irreducible.  Its first-coordinate image is
\(K+L=A\).  Since \(Z\to A\) is finite, this image has dimension \(n\)
and hence is all of \(Z\).  As \(\delta_Z\) is \(K\)-invariant,
\[
                      \delta_Z(E)
                       =\delta_Z(K\cdot E)
                       =\delta_Z(Z)=D .
\]
These are equalities of actual images; projectivity ensures that
no special point is lost by taking closures.
\end{proof}

\begin{proof}[Completion of the proof of Proposition~\ref{quot:smooth-base}]
Apply Lemmas~\ref{quot:ramification} and~\ref{quot:marked-divisor}
to the fixed \(D,K,Z,z_*\).  They provide a divisor \(L\subset A\),
a residual component \(\mathcal R'\subset\mathcal R\) containing
all of \(i(L)\), and a reduced prime divisor \(E\subset Z\) through
the translated point \(e_0\), with
\[
             \delta_Z(E)=D,\qquad
             \delta(\mathcal R')\not\subset D.
\]
Write \(x_0=s(e_0)=t(e_0)\in L\) and
\(\alpha_0=\nu(e_0)\).  The arrow \(\alpha_0\) is nonidentity and
has difference zero.  It is this translated arrow that composition
will retain.
Form the space of composable arrows
\[
                 \mathcal T
                  =\mathcal R'\times_{\,t,A,s} Z .
\]
The inclusion of all of \(i(L)\) in the actual, possibly nonnormal
component \(\mathcal R'\) gives a global morphism
\(i|_L:L\to\mathcal R'\): both \(L\) and \(\mathcal R'\) carry their
reduced structures.  Hence the map
\[
                    \iota:E\longrightarrow\mathcal T,\qquad
                    e\longmapsto \bigl(i(s(e)),e\bigr)
\]
is a closed embedding, the graph of this morphism over the closed
subscheme \(E\subset Z\).  Since \(E\) is irreducible, one irreducible
component \(\mathcal T_0\) contains its entire image.  Give
\(\mathcal T_0\) its reduced structure.  Since \(E\) is reduced,
\(\iota\) factors as a closed embedding into \(\mathcal T_0\);
in particular, this component contains the retained marked point.
Figure~\ref{quot:incidence} records this component choice and the
subsequent composition.

This same component dominates \emph{both} factors.
Indeed, \(\mathcal R'\times Z\) has pure dimension \(2n\).
The equations for \(\mathcal T\) are the pullback of the diagonal
in \(A\times A\), which is locally cut out by \(n\) equations.
The dimension inequality for an ideal with \(n\) generators implies
\(\dim\mathcal T_0\ge n\).
Both projections from \(\mathcal T\) to its factors are finite,
so \(\dim\mathcal T_0\le n\).
Its images in both irreducible \(n\)-dimensional factors are closed
of dimension \(n\), and thus are the whole factors.
The dimension estimate uses smoothness of \(A\) through the equations
of its diagonal.

Composition of the two labelled arrows gives
\(\mu:\mathcal T_0\to\mathcal R\).
Both \(\mathcal T_0\) and \(\mathcal R\) are finite over the outer source
coordinate in \(A\), so \(\mu\) is finite.  Its closed irreducible image
has dimension \(n\), since that source coordinate is surjective.
Consequently its image is an irreducible
component \(\mathcal R''\) of \(\mathcal R\).
On the embedded \(E\), composition with the identity is exactly
\(\nu(e)\), with its label.  Thus \(\mathcal R''\) contains the marked
nonidentity arrow \(\alpha_0=\nu(e_0)\) over difference zero, and its
difference image contains all of \(\delta_Z(E)=D\).

It also contains a point outside \(D\).  In fact, dominance of
\(\mathcal T_0\to\mathcal R'\) allows us to choose composable arrows
\[
                   z\longrightarrow x\longrightarrow y
\]
with \(x-z\notin D\), whereas \(y-x\in D\), because the second arrow
comes from \(Z\).  Since \(D\) is a subgroup,
\(y-z=(y-x)+(x-z)\notin D\).
The closed irreducible difference image of \(\mathcal R''\)
therefore contains \(D\) properly, and has dimension greater than
\(r\).  It still meets nonidentity inertia over zero, contradicting
the choice of \(r\).

\begin{figure}[tbp]
\centering
\begin{tikzpicture}[
  >=Stealth,
  object/.style={font=\large, inner sep=5pt, align=center},
  explanation/.style={font=\small, align=center},
  maplabel/.style={font=\small, fill=white, inner sep=2pt},
  retained/.style={draw=blue!55!black, line width=.8pt},
  projection/.style={draw=black!75, line width=.65pt}
]
  \node[object] (E) at (-4.5,0) {$E\ni e_0$};
  \node[object] (T) at (0,0) {$\mathcal T_0\ni\iota(e_0)$};
  \node[object] (Rpp) at (4.5,0) {$\mathcal R''\ni\alpha_0$};

  \node[explanation] at (0,1.05)
    {$\mathcal T_0$ is one irreducible component of
     $\mathcal R'\times_{t,A,s}Z$};

  \draw[retained,{Hooks[right]}->] (E.east) -- (T.west)
    node[midway,above,maplabel] {closed embedding $\iota$}
    node[midway,below,maplabel,font=\scriptsize]
      {$e\longmapsto\bigl(i(s(e)),e\bigr)$};
  \draw[retained,->] (T.east) -- (Rpp.west)
    node[midway,above,maplabel] {composition $\mu$}
    node[midway,below,maplabel,font=\scriptsize] {onto $\mathcal R''$};

  \node[object] (Rp) at (-2.05,-2.05) {$\mathcal R'$};
  \node[object] (Z) at (2.05,-2.05) {$Z$};
  \draw[projection,->] (T.south west) -- (Rp.north)
    node[midway,left=5pt,maplabel,font=\scriptsize,align=center]
      {finite,\\dominant};
  \draw[projection,->] (T.south east) -- (Z.north)
    node[midway,right=5pt,maplabel,font=\scriptsize,align=center]
      {finite,\\dominant};
  \node[explanation,font=\scriptsize,anchor=north] at (-2.05,-2.32)
    {actual relation component};
  \node[explanation,font=\scriptsize,anchor=north] at (2.05,-2.32)
    {normalization};

  \node[explanation,text=blue!55!black] at (0,-3.08)
    {$x_0=s(e_0),\qquad
      \mu\bigl(\iota(e_0)\bigr)=\alpha_0=\nu(e_0)\ne i(x_0),
      \qquad\delta(\alpha_0)=0$};

  \node[object,font=\normalsize] (source) at (-3.5,-4.15) {$z$};
  \node[object,font=\normalsize] (middle) at (0,-4.15) {$x$};
  \node[object,font=\normalsize] (target) at (3.5,-4.15) {$y$};
  \draw[projection,->] (source.east) -- (middle.west)
    node[midway,above,maplabel] {$\rho\in\mathcal R'$}
    node[midway,below,maplabel] {$x-z$};
  \draw[projection,->] (middle.east) -- (target.west)
    node[midway,above,maplabel] {$\nu(\zeta),\ \zeta\in Z$}
    node[midway,below,maplabel] {$y-x$};
  \node[explanation] at (0,-5.02)
    {$(\rho,\zeta)\in\mathcal T_0:\quad
      \mu(\rho,\zeta):z\longrightarrow y,$\\
      $y-z=(x-z)+(y-x)$};
\end{tikzpicture}
\caption{The component used for composition.  The whole divisor
\(E\), including \(e_0\), lies in one component \(\mathcal T_0\) whose
two finite projections are dominant.  Along \(E\) the first arrow is
the identity, so composition retains \(\nu(E)\) and the marked
nonidentity arrow \(\alpha_0\).  The bottom row shows how the two
differences add: dominance onto \(\mathcal R'\) supplies a difference
outside \(D\), while the second difference lies in \(D\).}
\label{quot:incidence}
\end{figure}

Thus every stabilizer is trivial.  A Deligne--Mumford stack
with trivial inertia is an algebraic space, and its coarse morphism
is then an isomorphism.  Therefore \(\mathfrak B=B\), and \(B\) is smooth.
\end{proof}
